% !TeX root = ../main.tex
\chapter{Funzioni}
\fonti{Lezioni 2, 4, 5, 6 e 7}

\section{Dominio, codominio e grafico}
% Fonti: L4 pp. 3--6; L5 p. 1; L6 p. 5.
\begin{definizione}
Una \textbf{funzione} $f:X\to Y$ associa a ogni $x\in X$ un unico elemento $f(x)\in Y$. $X$ è il \textbf{dominio}, $Y$ il \textbf{codominio}; l'\textbf{immagine} e il \textbf{grafico} sono
\[
\operatorname{Im}(f)=\{f(x)\mid x\in X\}\subseteq Y,
\qquad \operatorname{Gr}(f)=\{(x,f(x))\mid x\in X\}\subseteq X\times Y.
\]
Due funzioni sono uguali se hanno lo stesso dominio, lo stesso codominio e gli stessi valori in ogni punto del dominio.
\end{definizione}

\begin{esempio}[Una formula sui rappresentanti non basta]
La scrittura $f:\Q\to\Q$, $a/b\mapsto a+b$, non definisce una funzione: $1/2=2/4$, ma i due rappresentanti darebbero $3$ e $6$. Anche la regola $(a/b,c/d)\mapsto(a+c)/(bd)$ non definisce una funzione $\Q^2\to\Q$, come verificato nel capitolo precedente.
\end{esempio}

\begin{esempio}[Funzioni]
Consideriamo le seguenti funzioni:
\begin{enumerate}[label=\alph*)]
  \item $c:\C\to\C$, $c(z)=0$;
  \item $\operatorname{id}_{\R}:\R\to\R$, $x\mapsto x$;
  \item $s:\R_{\ge0}\to\R_{\ge0}$, $s(x)=x^2$;
  \item $a:T\to\R$, che associa a un triangolo non degenere del piano la sua area;
  \item $d:\N\times\N\to\Z$, $d(m,n)=m-n$.
\end{enumerate}
\end{esempio}

\begin{esempio}[Divisori primi]
Sia $P$ l’insieme dei primi positivi. La funzione $J:\N^*\to\mathcal P(P)$ associa a $n$ l’insieme $J(n)$ dei primi che dividono $n$. Per esempio $J(1)=\varnothing$, $J(4)=\{2\}$ e $J(6)=\{2,3\}$.
\end{esempio}

Una relazione è il grafico di una funzione da $X$ a $Y$ esattamente quando per ogni $x\in X$ esiste un unico $y\in Y$ a esso associato. Per un grafico nel piano, ogni retta verticale relativa a un punto del dominio deve dunque incontrare il grafico in un solo punto. La circonferenza unitaria non è il grafico di una funzione da $[-1,1]$ a $\R$, perché alla stessa ascissa interna corrispondono due ordinate.

\begin{definizione}
Una funzione $f:X\to Y$ è
\begin{itemize}
  \item \textbf{iniettiva} se $f(x)=f(x')$ implica $x=x'$ per ogni $x,x'\in X$;
  \item \textbf{suriettiva} se per ogni $y\in Y$ esiste $x\in X$ con $f(x)=y$, cioè se $\operatorname{Im}(f)=Y$.
\end{itemize}
\end{definizione}

\begin{esercizio}[Primi esempi]
Discutere iniettività e suriettività delle cinque funzioni elencate sopra.
\end{esercizio}
\begin{svolgimento}
\textbf{Costante.} $c(0)=c(1)$, quindi $c$ non è iniettiva. La sua immagine è $\{0\}$, perciò non è suriettiva.

\textbf{Identità.} L'uguaglianza delle immagini è già l'uguaglianza dei punti e ogni $y$ è immagine di se stesso: è iniettiva e suriettiva. In generale $\operatorname{id}_X(x)=x$ definisce l'identità di qualsiasi insieme $X$.

\textbf{Quadrato sui non negativi.} Se $x^2=y^2$ e $x,y\ge0$, da $(x-y)(x+y)=0$ segue $x=y$. Ogni $t\ge0$ ha preimmagine $\sqrt t$. È quindi iniettiva e suriettiva. Anche la restrizione del quadrato a $\R_{\le0}\to\R_{\ge0}$ è iniettiva e suriettiva: da due quadrati uguali di numeri non positivi segue l'uguaglianza dei numeri, e $t\ge0$ è immagine di $-\sqrt t$. Con dominio $\R$ si perderebbe l'iniettività, poiché $s(1)=s(-1)$; con codominio $\R$ si perderebbe la suriettività.

\textbf{Area.} Triangoli distinti possono avere la stessa area, per esempio due triangoli ottenuti uno dall'altro per traslazione. L'immagine è $(0,+\infty)$: un triangolo di base $2$ e altezza $t>0$ ha area $t$. La funzione non è dunque né iniettiva né suriettiva su $\R$.

\textbf{Differenza.} Non è iniettiva: $d(0,0)=d(1,1)=0$. È suriettiva: se $z\ge0$, $d(z,0)=z$; se $z<0$, $d(0,-z)=z$.
\end{svolgimento}

% A–L L4 p. 5; svolgimento confermato dal ritorno M–Z L7 pp. 1--2.
\begin{esercizio}[Una funzione a tratti sugli interi]
Studiare $f:\Z\to\Z$ definita da
\[
f(n)=\begin{cases}2n-3&n\text{ pari},\\n+1&n\text{ dispari}.\end{cases}
\]
\end{esercizio}
\begin{svolgimento}
Scrivendo $n=2k$, si ottiene $f(n)=4k-3$, che percorre gli interi della forma $4k+1$, con $k\in\Z$. Per $n=2k+1$ si ottiene $f(n)=2k+2$, che percorre tutti gli interi pari. Le due immagini sono disgiunte e ciascuna regola è iniettiva sul proprio ramo: quindi $f$ è iniettiva. Non è suriettiva, perché mancano gli interi della forma $4k+3$, per esempio $3$.
\end{svolgimento}

\begin{esercizio}[Una funzione del piano]
Studiare $F:\R^2\to\R^2$, $F(x,y)=(y^3,x(y^2+1))$.
\end{esercizio}
\begin{svolgimento}
% Supporto M–Z L5 pp. 2--3: verifiche separate, senza cambiare il punto A–L.
\textbf{Iniettività.} Se $F(x,y)=F(x',y')$, dalla prima coordinata segue $y^3=(y')^3$, dunque $y=y'$. Dalla seconda, $x(y^2+1)=x'(y^2+1)$; poiché $y^2+1>0$, segue $x=x'$.

\textbf{Suriettività.} Per ogni $(u,v)\in\R^2$ poniamo
\[
y=\sqrt[3]u,\qquad x=\frac{v}{1+(\sqrt[3]u)^2}.
\]
Il denominatore è positivo. Sostituendo si ottiene $F(x,y)=(u,v)$, quindi ogni punto del codominio è raggiunto.
\end{svolgimento}

\section{Fibre e proiezione sul quoziente}
% Fonti: L2 p. 5; L4 pp. 6--7; L5 p. 1; L6 pp. 1--2.
% M–Z L5 pp. 4--5 conferma il calcolo delle fibre. Estensione a 0: ripresa A–L L6.
\begin{definizione}
La \textbf{fibra} di $f:X\to Y$ sopra $y\in Y$ è $\{x\in X\mid f(x)=y\}$. Per $x\in X$ la classe dell'equivalenza $x\sim_f x'\iff f(x)=f(x')$ è la fibra sopra $f(x)$:
\[
[x]_{\sim_f}=\{x'\in X\mid f(x')=f(x)\}.
\]
Le fibre non vuote formano una partizione di $X$. Le fibre sopra punti fuori dall'immagine sono vuote e non sono classi.
\end{definizione}

\begin{schema}
\begin{tikzpicture}[>=Stealth]
  \draw (0,0) ellipse (1.4 and 1.15);\draw (5,0) ellipse (1.1 and 1.15);
  \node at (-1.3,1.2) {$X$};\node at (6,1.2) {$Y$};
  \draw[dashed,coloreschema] (-.2,.25) ellipse (.8 and .65);
  \node (a) at (-.35,.6) {$x_1$};\node (b) at (-.15,-.1) {$x_2$};\node (c) at (0,-.8) {$x_3$};
  \node (u) at (4.8,.6) {$y_0$};\node (v) at (5,-.15) {$y_1$};\node at (5,-.85) {$y_2$};
  \draw[->,coloreschema] (a)--(u);\draw[->,coloreschema] (b)--(u);\draw[->,coloreschema] (c)--(v);
  \node at (2.5,-1.65) {Una fibra può contenere più punti; un valore può non avere preimmagini.};
\end{tikzpicture}
\end{schema}

% Integrazione M–Z L5 p. 4.
\begin{proposizione}
La funzione $f:X\to Y$ è iniettiva se e solo se l'equivalenza $x\sim_f x'\iff f(x)=f(x')$ è l'uguaglianza su $X$.
\end{proposizione}
\begin{proof}
Per iniettività, $x\sim_f x'$ implica $x=x'$; il viceversa è sempre vero. Se invece $\sim_f$ è l'uguaglianza, da $f(x)=f(x')$ segue $x\sim_f x'$, quindi $x=x'$ e $f$ è iniettiva.
\end{proof}

\begin{esempio}
Per $f:\R\to\R$, $f(x)=x^2$, si ha $[2]_{\sim_f}=[-2]_{\sim_f}=\{2,-2\}$, mentre $[0]_{\sim_f}=\{0\}$. La fibra sopra $-1$ è vuota.
\end{esempio}

\begin{esercizio}[Insiemi di divisori primi]
Sia $P$ l'insieme dei primi positivi e sia
\[
J:\N\to\mathcal P(P),\qquad J(n)=\{p\in P\mid p\mid n\}.
\]
Determinare immagini, iniettività, suriettività e le classi di $0,1,4,5,6,12,168$, considerando anche la restrizione a $\N^*$.
\end{esercizio}
\begin{svolgimento}
Si ha
\[
\begin{gathered}
J(0)=P,\quad J(1)=\varnothing,\quad J(4)=\{2\},\quad J(12)=\{2,3\},\\
J(30)=\{2,3,5\},\qquad J(7)=\{7\}.
\end{gathered}
\]
Ogni primo divide $0$ perché $0=p\cdot0$, mentre nessun primo divide $1$. La funzione non è iniettiva: $J(2)=J(4)$.

Per $n>0$ l'insieme $J(n)$ è finito, perché ogni suo elemento è compreso tra $2$ e $n$. Viceversa ogni sottoinsieme finito $\{p_1,\ldots,p_r\}$ è immagine del prodotto $p_1\cdots p_r$, ponendo il prodotto vuoto uguale a $1$. Pertanto
\[
\begin{aligned}
\operatorname{Im}(J|_{\N^*})&=\{A\subseteq P\mid A\text{ finito}\},\\
\operatorname{Im}(J)&=\{A\subseteq P\mid A\text{ finito}\}\cup\{P\}.
\end{aligned}
\]
In questo esempio si usano le proprietà aritmetiche usuali dei primi: sono infiniti e ogni intero positivo si fattorizza in primi.\footnote{Questi fatti sono qui assunti; la loro trattazione e le prove sono nel capitolo~\ref{cap:interi}.} In particolare $P\setminus\{2\}$ è infinito e diverso da $P$, quindi non appartiene a nessuna delle due immagini: entrambe le funzioni non sono suriettive. Si ottengono le fibre
\[
[5]_{\sim_J}=\{5^k\mid k\ge1\},\qquad
[6]_{\sim_J}=\{2^a3^b\mid a,b\ge1\},\qquad [0]_{\sim_J}=\{0\}.
\]
Inoltre, poiché $168=2^3\cdot3\cdot7$,
\[
[1]_{\sim_J}=\{1\},\quad [4]_{\sim_J}=\{2^k\mid k\ge1\},\quad
[12]_{\sim_J}=[6]_{\sim_J},
\]
\[
[168]_{\sim_J}=\{2^a3^b7^c\mid a,b,c\ge1\}.
\]
Gli esponenti devono essere positivi perché ciascuno dei primi indicati deve dividere il numero; non possono comparire altri primi.
% M–Z L5 pp. 4--5, compreso l'esercizio su [168].
\end{svolgimento}

\begin{definizione}
Una funzione è \textbf{biettiva}, o una \textbf{biezione}, se è iniettiva e suriettiva.
\end{definizione}

\begin{proposizione}[Proiezione canonica]
Per ogni equivalenza $\sim$ su $X$, la funzione $\pi:X\to X/{\sim}$, $x\mapsto[x]$, è suriettiva e le sue fibre sono le classi di $\sim$.
\end{proposizione}
\begin{proof}
Ogni elemento del codominio ha la forma $[x]=\pi(x)$ per definizione del quoziente. Inoltre $\pi(x')=\pi(x)$ equivale a $[x']=[x]$, dunque a $x'\sim x$. La fibra sopra $[x]$ è pertanto $[x]$ stesso, considerato ora come sottoinsieme del dominio.
\end{proof}

\section{Immagine diretta e controimmagine}
% Fonte: L5 p. 2 e p. 4.
\begin{definizione}
A una funzione $f:X\to Y$ si associano
\[
\begin{aligned}
f_*:\mathcal P(X)&\to\mathcal P(Y),&f_*(A)&=\{f(a)\mid a\in A\},\\
f^*:\mathcal P(Y)&\to\mathcal P(X),&f^*(B)&=\{x\in X\mid f(x)\in B\}.
\end{aligned}
\]
Si chiamano rispettivamente \textbf{immagine diretta} e \textbf{controimmagine}. Si scrive spesso $f(A)$ per $f_*(A)$. Per un punto $y$, $f^*(\{y\})$ è la fibra sopra $y$.
\end{definizione}

\begin{esempio}[Il quadrato]
Per $f:\R\to\R$, $f(x)=x^2$, e
\[
A=[-1,2],\quad A'=[-2,0],\quad B=[0,4],\quad B'=(-\infty,4],
\]
si ha $f_*(A)=f_*(A')=[0,4]$ e $f^*(B)=f^*(B')=[-2,2]$. Infatti i quadrati dei punti di entrambi gli intervalli percorrono $[0,4]$; la condizione $x^2\le4$ equivale a $-2\le x\le2$ e il vincolo $x^2\ge0$ è automatico.

In particolare $f^*(f_*(A))=[-2,2]\ne A$. La controimmagine non annulla in generale l'immagine diretta: recupera tutte le fibre che incontrano $A$. Inoltre $f^*((-\infty,0))=\varnothing$, perché un quadrato non è negativo.
\end{esempio}

\begin{schema}
\begin{tikzpicture}[scale=.85]
  \draw[asse] (-2.8,0)--(2.8,0) node[right] {$x$};\draw[asse] (0,-.3)--(0,4.9) node[above] {$y$};
  \draw[coloreschema,thick,domain=-2.1:2.1,samples=60] plot (\x,{\x*\x});
  \draw[densely dashed] (-2,4)--(2,4);\draw[densely dashed] (-2,0)--(-2,4);\draw[densely dashed] (2,0)--(2,4);
  \node[below] at (-2,0) {$-2$};\node[below] at (-1,0) {$-1$};\node[below] at (2,0) {$2$};\node[above left] at (0,4) {$4$};
  \draw[coloreschema,line width=2pt] (-1,-.35)--(2,-.35);\node[right] at (2.2,-.35) {$A$};
\end{tikzpicture}
\end{schema}

% Integrazione M–Z L5 pp. 7--8: proprietà generali oltre all'esempio del quadrato.
\begin{proposizione}
Per $f:X\to Y$, $A\subseteq X$ e $B\subseteq Y$,
\[
A\subseteq f^*(f_*(A)),\qquad f_*(f^*(B))=B\cap\operatorname{Im}(f)\subseteq B.
\]
La prima inclusione è un'uguaglianza per ogni $A$ se e solo se $f$ è iniettiva. La seconda è un'uguaglianza con $B$ per ogni $B$ se e solo se $f$ è suriettiva.
\end{proposizione}
\begin{proof}
\textbf{Inclusioni.} Se $x\in A$, allora $f(x)\in f_*(A)$, dunque $x\in f^*(f_*(A))$. Un punto $y$ appartiene a $f_*(f^*(B))$ esattamente quando è immagine di un $x$ con $f(x)\in B$: ciò equivale a $y\in B\cap\operatorname{Im}(f)$.

\textbf{Iniettività.} Se $f$ è iniettiva e $x\in f^*(f_*(A))$, esiste $a\in A$ con $f(x)=f(a)$; quindi $x=a\in A$. Viceversa, supponiamo l'uguaglianza per tutti i sottoinsiemi e siano $f(x)=f(x')$. Prendendo $A=\{x\}$, si ha $x'\in f^*(f_*(A))=A$, dunque $x'=x$.

\textbf{Suriettività.} Se $\operatorname{Im}(f)=Y$, la formula dà $f_*(f^*(B))=B$. Viceversa, applicando l'uguaglianza a $B=Y$, si ottiene $\operatorname{Im}(f)=f_*(X)=Y$.
\end{proof}

\section{Composizione}
% Fonti: L5 pp. 3--6.
\begin{definizione}
Per $f:X\to Y$ e $g:Y\to Z$, la \textbf{composizione} è
\[
g\circ f:X\to Z,\qquad (g\circ f)(x)=g(f(x)).
\]
La compatibilità fra codominio della prima funzione e dominio della seconda è parte della definizione. Per esempio una funzione da $\{1,2,3\}$ a $\N$, seguita dalla funzione $\N\to\R$, $n\mapsto n$, è una funzione da $\{1,2,3\}$ a $\R$.
\end{definizione}

\begin{proposizione}[Associatività]
Se $f:X\to Y$, $g:Y\to Z$, $h:Z\to T$, allora
\[
h\circ(g\circ f)=(h\circ g)\circ f.
\]
\end{proposizione}
\begin{proof}
Entrambe le composizioni hanno dominio $X$, codominio $T$ e valore $h(g(f(x)))$ in ogni $x\in X$.
\end{proof}

\begin{proposizione}
Per $f:X\to Y$ e $g:Y\to Z$ valgono
\[
(g\circ f)_*=g_*\circ f_*:\mathcal P(X)\to\mathcal P(Z),
\]
\[
(g\circ f)^*=f^*\circ g^*:\mathcal P(Z)\to\mathcal P(X).
\]
L'immagine diretta conserva l'ordine di composizione; la controimmagine lo inverte.
\end{proposizione}
\begin{proof}
Per $A\subseteq X$, entrambi i membri della prima identità sono $\{g(f(a))\mid a\in A\}$. Per $C\subseteq Z$,
\[
x\in(g\circ f)^*(C)\iff g(f(x))\in C
\iff f(x)\in g^*(C)\iff x\in f^*(g^*(C)).
\]
Questo prova la seconda identità per ogni sottoinsieme del codominio.
\end{proof}

\begin{proposizione}[Iniettività e suriettività delle composizioni]\label{prop:composte}
Date $f:X\to Y$ e $g:Y\to Z$:
\begin{enumerate}[label=\alph*)]
  \item se $f$ e $g$ sono iniettive, anche $g\circ f$ è iniettiva;
  \item se $f$ e $g$ sono suriettive, anche $g\circ f$ è suriettiva;
  \item se $g\circ f$ è iniettiva, $f$ è iniettiva;
  \item se $g\circ f$ è suriettiva, $g$ è suriettiva.
\end{enumerate}
\end{proposizione}
\begin{proof}
\textbf{a).} Se $g(f(x))=g(f(x'))$, l'iniettività di $g$ dà $f(x)=f(x')$; quella di $f$ dà $x=x'$.

\textbf{b).} Dato $z\in Z$, la suriettività di $g$ fornisce $y\in Y$ con $g(y)=z$; quella di $f$ fornisce $x\in X$ con $f(x)=y$. Quindi $(g\circ f)(x)=z$.

\textbf{c).} Da $f(x)=f(x')$ segue $g(f(x))=g(f(x'))$, e l'iniettività della composta impone $x=x'$.

\textbf{d).} Per $z\in Z$, esiste $x\in X$ con $g(f(x))=z$. L'elemento $y=f(x)\in Y$ è una preimmagine di $z$ per $g$.
\end{proof}

\begin{esempio}[Limiti delle implicazioni]
I conversi delle ultime due affermazioni sono falsi:
\begin{itemize}
  \item $f(x)=x^3$ e $g(y)=y^2$, entrambe su $\R$, hanno $f$ iniettiva ma $g\circ f(x)=x^6$ non iniettiva;
  \item $f:\R\to\R$, $f(x)=x^2$, e $g:\R\to\R_{\ge0}$, $g(y)=|y|$, hanno composta suriettiva benché $f$ non sia suriettiva su $\R$;
  \item l'inclusione $f:\R_{\ge0}\hookrightarrow\R$, seguita da $g:\R\to\R$, $g(y)=y^2$, dà una composta iniettiva benché $g$ non sia iniettiva;
  \item l'inclusione $f:[0,1]\hookrightarrow\R$, seguita da $g:\R\to\R_{\ge0}$, $g(y)=y^2$, dà una composta con immagine $[0,1]$, benché $g$ sia suriettiva.
\end{itemize}
Ogni affermazione dipende anche dal dominio e dal codominio dichiarati.
\end{esempio}

\section{Inverse a sinistra, a destra e bilatere}
% Fonti: L5 p. 7; L6 pp. 1--5.
Per ogni insieme $X$, l'identità è $\operatorname{id}_X:X\to X$, $x\mapsto x$. Per $f:X\to Y$ si ha
\[
\operatorname{id}_Y\circ f=f=f\circ\operatorname{id}_X,
\]
poiché entrambi i membri hanno valore $f(x)$ in ogni $x\in X$.

\begin{definizione}
Per $f:X\to Y$, una funzione $g:Y\to X$ è un'\textbf{inversa a sinistra} se $g\circ f=\operatorname{id}_X$; è un'\textbf{inversa a destra} se $f\circ g=\operatorname{id}_Y$. Una funzione $f$ è \textbf{invertibile} se possiede un'inversa bilatera, indicata con $f^{-1}$.
\end{definizione}

\begin{proposizione}[Unicità dell'inversa bilatera]
Se $g\circ f=\operatorname{id}_X$ e $f\circ h=\operatorname{id}_Y$, allora $g=h$. In particolare, se esistono un'inversa a sinistra e una a destra, esse coincidono.
\end{proposizione}
\begin{proof}
Usando associatività e identità,
\[
g=g\circ\operatorname{id}_Y=g\circ(f\circ h)
=(g\circ f)\circ h=\operatorname{id}_X\circ h=h.
\]
\end{proof}

\begin{esempio}[Inverse a un solo lato non uniche]
Sia $f:\R\to\R_{\ge0}$, $f(x)=x^2$. Le due funzioni $g_+(y)=\sqrt y$ e $g_-(y)=-\sqrt y$ sono inverse a destra di $f$, poiché $f(g_\pm(y))=y$. Ma $g_+\circ f(x)=|x|$, che non è $\operatorname{id}_{\R}$: per $x=-1$ si ottiene $1$ anziché $-1$. La funzione $f$ non è invertibile, perché non è iniettiva.
\end{esempio}

\begin{teorema}[Criteri di esistenza delle inverse]\label{teo:inverse}
Sia $f:X\to Y$, con $X\ne\varnothing$.
\begin{enumerate}[label=\alph*)]
  \item $f$ è iniettiva se e solo se ammette un'inversa a sinistra.
  \item $f$ è suriettiva se e solo se ammette un'inversa a destra, nel quadro con assioma della scelta.
  \item $f$ è biettiva se e solo se è invertibile.
\end{enumerate}
\end{teorema}
\begin{proof}
\textbf{a), costruzione.} Se $f$ è iniettiva, ogni $y\in\operatorname{Im}(f)$ ha un'unica preimmagine $x_y$. Fissato $x_0\in X$, poniamo
\[
g(y)=\begin{cases}x_y&y\in\operatorname{Im}(f),\\x_0&y\notin\operatorname{Im}(f).\end{cases}
\]
È una funzione $Y\to X$ e $g(f(x))=x$ per ogni $x$. Viceversa, se $g\circ f=\operatorname{id}_X$, la composta è iniettiva, quindi $f$ è iniettiva per la proposizione~\ref{prop:composte}. Ogni tale $g$ è inoltre suriettiva.

\textbf{b), costruzione.} Se $f$ è suriettiva, ogni fibra sopra $y\in Y$ è non vuota. Scegliendo simultaneamente un rappresentante $x_y$ per ogni fibra si definisce $g(y)=x_y$; allora $f(g(y))=y$. Questa scelta per una famiglia arbitraria di fibre è il punto in cui si usa l'assioma della scelta. Viceversa $f\circ g=\operatorname{id}_Y$ implica che $f$ è suriettiva. Ogni tale $g$ è iniettiva, perché una sua composta a sinistra è l'identità.

\textbf{c).} Per una biezione ogni $y$ ha un'unica preimmagine: la funzione che gliela associa è inversa da entrambi i lati, senza scelte arbitrarie. Se invece esiste un'inversa bilatera, le due identità e i versi già dimostrati impongono sia iniettività sia suriettività.
\end{proof}

\begin{osservazione}[Insiemi vuoti]
La condizione $X\ne\varnothing$ nel primo criterio consente di scegliere il valore $x_0$ fuori dall'immagine. La funzione vuota $\varnothing\to Y$ è iniettiva, ma per $Y\ne\varnothing$ non esiste alcuna funzione $Y\to\varnothing$. La funzione $\varnothing\to\varnothing$ è invece biettiva e inversa di se stessa.
\end{osservazione}

\begin{corollario}
Per insiemi non vuoti $X,Y$, esiste un'iniezione $X\to Y$ se e solo se esiste una suriezione $Y\to X$, nel quadro con assioma della scelta.
\end{corollario}
\begin{proof}
Un'iniezione ammette un'inversa a sinistra, che è suriettiva. Una suriezione ammette un'inversa a destra, che è iniettiva. Le funzioni ottenute hanno nei due casi il verso richiesto.
\end{proof}

Si usano $X\hookrightarrow Y$ per un'iniezione, $X\twoheadrightarrow Y$ per una suriezione e $X\xrightarrow{\sim}Y$ per una biezione. Se $X\subseteq Y$, la funzione $i:X\hookrightarrow Y$, $i(x)=x$, si chiama \textbf{inclusione}.

% Integrazione M–Z L6 p. 6: esercizio sulle inverse destre del seno.
\begin{esercizio}[Inverse destre del seno]
Trovare due inverse a destra distinte di $s:\R\to[-1,1]$, $s(x)=\sin x$. Esiste un'inversa a sinistra?
\end{esercizio}
\begin{svolgimento}
Si usa la proprietà analitica usuale che $\sin$ ristretto a $[-\pi/2,\pi/2]$ è una biezione su $[-1,1]$. La sua inversa è $\arcsin:[-1,1]\to[-\pi/2,\pi/2]$. Poniamo
\[
g_1(y)=\arcsin y,\qquad g_2(y)=\pi-\arcsin y.
\]
Per $y\in[-1,1]$, $\sin(g_1(y))=y$ e, usando $\sin(\pi-t)=\sin t$, anche $\sin(g_2(y))=y$. Sono distinte, perché $g_1(0)=0$ e $g_2(0)=\pi$. Non esiste un'inversa a sinistra: $s(0)=s(\pi)=0$ rende $s$ non iniettiva.
\end{svolgimento}

\section{Insiemi di funzioni e funzioni caratteristiche}
% Fonti: L6 pp. 6--7; L7 pp. 1--2.
\begin{definizione}
Per insiemi $X,\Omega$ si scrive
\[
\Omega^X=\{\varphi\mid \varphi:X\to\Omega\text{ è una funzione}\}.
\]
Una funzione $\{1,\ldots,n\}\to\Omega$ equivale a una $n$-upla di elementi di $\Omega$, mediante $\varphi\mapsto(\varphi(1),\ldots,\varphi(n))$. Perciò $\R^4$ si identifica con $\R^{\{1,2,3,4\}}$; $\R^\N$ è l'insieme delle successioni $(x_0,x_1,\ldots)$ reali. Più in generale, una funzione $I\to\Omega$ si identifica con una famiglia $(x_i)_{i\in I}$ di elementi di $\Omega$, ponendo $x_i=\varphi(i)$: l'indice determina un unico elemento, mentre indici diversi possono avere lo stesso valore. % M–Z L7 pp. 5--7.

\end{definizione}

\begin{definizione}
Per $A\subseteq X$, la \textbf{funzione caratteristica} di $A$ è
\[
\chi_A:X\to\{0,1\},\qquad
\chi_A(x)=\begin{cases}1&x\in A,\\0&x\notin A.\end{cases}
\]
\end{definizione}

\begin{teorema}\label{teo:caratteristiche}
Le applicazioni
\[
\begin{aligned}
\alpha:\{0,1\}^X&\to\mathcal P(X),&\alpha(\varphi)&=\varphi^*(\{1\}),\\
\beta:\mathcal P(X)&\to\{0,1\}^X,&\beta(A)&=\chi_A
\end{aligned}
\]
sono biezioni inverse.
\end{teorema}
\begin{proof}
\passomath{\alpha\circ\beta=\operatorname{id}_{\mathcal P(X)}} Per ogni $A\subseteq X$,
\[
\alpha(\beta(A))=\{x\in X\mid\chi_A(x)=1\}=A.
\]
\passomath{\beta\circ\alpha=\operatorname{id}_{\{0,1\}^X}} Per $\varphi:X\to\{0,1\}$, la funzione caratteristica di $\varphi^*(\{1\})$ vale $1$ esattamente nei punti in cui $\varphi$ vale $1$. Negli altri punti entrambe valgono $0$, perché gli unici valori ammessi sono $0$ e $1$. Quindi $\chi_{\varphi^*(\{1\})}=\varphi$.

Le due identità provano che le applicazioni sono inverse e biettive. L'argomento vale anche per $X=\varnothing$: in quel caso entrambi gli insiemi hanno un unico elemento.
\end{proof}
